\documentclass[10pt,a4paper]{amsart}
\usepackage[margin=19mm]{geometry}
\usepackage{amsmath,amssymb,mathtools}
\usepackage[scr=rsfs,cal=euler]{mathalfa}
\usepackage{xfrac}
\usepackage[unicode,hidelinks,bookmarks=false]{hyperref}
\newcommand{\ie}{i.e.\ }
\newcommand{\cf}{cf.\ }
\newcommand{\resp}{respectively\ }
\newcommand{\Subsection}{\subsection}
\providecommand{\nobreakdash}{\nobreak\hbox{-}}
\setlength{\emergencystretch}{2em}
\multlinegap0pt
\usepackage{bm}
\usepackage{bbm}
\usepackage{dsfont}
\usepackage{xspace}
\usepackage{cancel}
\usepackage{mathtools}
\usepackage{amsthm}
\makeatletter
\@ifundefined{lemma}{\newtheorem{lemma}{Lemma}}{}
\makeatother
\newcommand{\R}{\mathbb{R}}
\newcommand{\Sph}{\mathbb S}
\newcommand{\abs}[1]{\left|#1\right|}
\newcommand{\ddsigma}{\,d\sigma}
\newcommand{\norm}[1]{\left\|#1\right\|}
\title[Sharp Sobolev Sandwich and Approximation Rates of Radon-Doma]{Sharp Sobolev Sandwich and Approximation Rates of Radon-Domain $L^p$ Ridge Integral Spaces for ReLU$^k$ Networks}
\author{Juncai He, Zitong Tian}
\date{Source: arXiv:2606.24795v2 (CC BY 4.0)}
\begin{document}
\maketitle

\noindent\emph{Source and licence.} Sharp Sobolev Sandwich and Approximation Rates of Radon-Domain $L^p$ Ridge Integral Spaces for ReLU$^k$ Networks. Version arXiv:2606.24795v2 (CC BY 4.0), distributed under the Creative Commons Attribution 4.0 International licence, as recorded in the arXiv licence metadata for this exact version and shown on the abstract page https://arxiv.org/abs/2606.24795v2. Full rights notice: licenses/arxiv-2606.24795-CC-BY-4.0.txt.\par\smallskip

\noindent\textit{Editorial scope.} Counted unit: the Lemma whose source label is \texttt{lem:localized-dual-radon} in the article (source0.tex lines 877--886, source label \texttt{lem:localized-dual-radon}), delivered together with the article's own complete original proof of it (source0.tex lines 888--954, one \texttt{proof} environment of 67 lines). No mathematics is written, repaired or re-derived: statement and proof are the article's own text line for line. Required context retained uncounted: eq:antipodal (lines 565--569). Every definition, display and label of those lines is retained unchanged. Omitted from this package and not redistributed: 1--564, 570--876, 887--887, 955--3174. Nothing inside a retained excerpt is omitted or altered: every retained body is delivered exactly as the article's lines are written, so no omission receipt of any kind accompanies this package. Citations: the retained excerpts carry no citation command at all, so no bibliography entry of the article is transcribed and no reference is redistributed. Reference adaptation: none was needed and none was made; every reference inside the delivered unit resolves inside it, so no unresolved reference or citation is produced. Printed slips kept literal: no printed word, symbol, hypothesis or number of the counted statement or of its proof was altered. Layout and class: the article's own class is \texttt{article} with packages including hyperref, geometry, color, amssymb, bm, amsmath, amsfonts, amsthm, mathrsfs, appendix, tikz, natbib, array, booktabs; the wrapper is the lane's pinned \texttt{amsart} editorial wrapper at 10pt on A4, which restates the theorem-like environments the retained text needs and adds the article's own macro definitions that the retained text needs (\texttt{\textbackslash R}, \texttt{\textbackslash Sph}, \texttt{\textbackslash abs}, \texttt{\textbackslash ddsigma}, \texttt{\textbackslash norm}), with extra wrapper packages (bm, bbm, dsfont, xspace, cancel, mathtools). The delivered numbering is the unit's own. Assets: no retained span contains a figure environment, a graphics-inclusion command, a PDF-page inclusion, a table or a TikZ picture; no image file is copied into this package, the package declares no asset dependency and no image file is redistributed with it. Licence: version arXiv:2606.24795v2 of this article is distributed under the Creative Commons Attribution 4.0 International licence, as recorded in the arXiv licence metadata for this exact version and shown on the abstract page https://arxiv.org/abs/2606.24795v2. A full rights notice is delivered at licenses/arxiv-2606.24795-CC-BY-4.0.txt. Characters: every retained excerpt is pure ASCII, so the delivered engine, XeLaTeX with its default Latin Modern text font, reports no missing character.
\begin{align}
    \int_{\Sph^{d-1}}\!\!\int_\R \abs{\omega}^{d-1}h(w,\omega)\,d\omega\,\ddsigma(w)
    =2\int_{\R^d} h\big(\tfrac{\xi}{\abs{\xi}},\abs{\xi}\big)\,d\xi .
    \label{eq:antipodal}
\end{align}

\begin{lemma}
\label{lem:localized-dual-radon}
Let $\eta\in C_c^\infty(\R^d)$, $\psi\in C_c^\infty(\Sph^{d-1}\times\R)$, and
$r\ge0$. Then for every $\Phi\in L^2(\Sph^{d-1};H^r(\R))$,
\begin{align}
    \norm{\eta\,\mathcal R^*(\psi\Phi)}_{H^{r+(d-1)/2}(\R^d)}
    \le C_{\eta,\psi,r}\,\norm{\Phi}_{L^2(\Sph^{d-1};H^r(\R))}.
    \label{eq:localized-dual-radon-estimate}
\end{align}
\end{lemma}

\begin{proof}
By density it suffices to treat $\Phi\in C_c^\infty(\Sph^{d-1}\times\R)$. For each
fixed $w$, multiplication by $\psi(w,\cdot)\in C_c^\infty(\R)$ is bounded on
$H^r(\R)$, with operator norm controlled by finitely many derivatives of $\psi$
uniformly in $w$; hence multiplication by $\psi$ is bounded on
$L^2(\Sph^{d-1};H^r(\R))$. We may therefore prove
\eqref{eq:localized-dual-radon-estimate} with right-hand side
$\norm{\psi\Phi}_{L^2(\Sph^{d-1};H^r(\R))}$ and absorb the multiplier constant at
the end. Write $\Psi:=\psi\Phi$ and $G:=\mathcal F_b\Psi$. Fourier inversion in the
second variable gives
\begin{align*}
    \widehat{\eta\,\mathcal R^*\Psi}(\xi)
    =\frac1{2\pi}\int_{\Sph^{d-1}}\int_\R
    \widehat\eta(\xi-\omega w)\,G(w,\omega)\,d\omega\,d\sigma(w).
\end{align*}
Fix $\xi\in\R^d$ and set $\langle\omega\rangle:=(1+|\omega|^2)^{1/2}$. Factoring the
integrand as
$|\widehat\eta(\xi-\omega w)|^{1/2}\langle\omega\rangle^{-r}
\cdot|\widehat\eta(\xi-\omega w)|^{1/2}\langle\omega\rangle^{r}|G(w,\omega)|$
and applying Cauchy--Schwarz in $(w,\omega)$,
\begin{align}
    |\widehat{\eta\,\mathcal R^*\Psi}(\xi)|^2
    \le\frac1{(2\pi)^2}\,K(\xi)\,J(\xi),
    \label{eq:CS-split}
\end{align}
where
\begin{align*}
    K(\xi):&=\int_{\Sph^{d-1}}\int_\R|\widehat\eta(\xi-\omega w)|
\langle\omega\rangle^{-2r}\,d\omega\,d\sigma(w), \\
    J(\xi):&=\int_{\Sph^{d-1}}\int_\R|\widehat\eta(\xi-\omega w)|
\langle\omega\rangle^{2r}|G(w,\omega)|^2\,d\omega\,d\sigma(w).
\end{align*}
Denote $h(\zeta):=\langle\zeta\rangle^{-2r}|\zeta|^{-(d-1)}$, so that
$K=2\,|\widehat\eta|*h$ by polar coordinates~\eqref{eq:antipodal}. Since
$h\notin L^1(\R^d)$ when $r\le\tfrac12$, we read off the decay of $K$ from the
convolution itself rather than from $\norm{h}_{L^1}$. By Peetre's inequality
$\langle\xi\rangle^{2r+d-1}\le C\langle\xi-\zeta\rangle^{2r+d-1}
\langle\zeta\rangle^{2r+d-1}$ together with
$\langle\zeta\rangle^{2r+d-1}h(\zeta)=\langle\zeta\rangle^{d-1}|\zeta|^{-(d-1)}
\le C\bigl(1+|\zeta|^{-(d-1)}\bigr)$, the factor
$\langle\,\cdot\,\rangle^{2r+d-1}|\widehat\eta|\in L^1\cap L^\infty(\R^d)$
convolved against $1+|\zeta|^{-(d-1)}\in L^\infty+L^1_{\mathrm{loc}}$ is bounded
uniformly in $\xi$. Hence $K(\xi)\le C_0\langle\xi\rangle^{-(2r+d-1)}$ for all
$r\ge0$, with $C_0=C_0(\eta,r,d)$. Inserting this into
\eqref{eq:CS-split} gives
$\langle\xi\rangle^{2r+d-1}|\widehat{\eta\,\mathcal R^*\Psi}(\xi)|^2
\le\tfrac{C_0}{(2\pi)^2}J(\xi)$. Integrating in $\xi$, and applying Tonelli's
theorem together with the translation invariance
$\int_{\R^d}|\widehat\eta(\xi-\omega w)|\,d\xi=\norm{\widehat\eta}_{L^1(\R^d)}$
(independent of $(w,\omega)$), we obtain
\begin{align}
    \int_{\R^d}\langle\xi\rangle^{2r+d-1}
    |\widehat{\eta\,\mathcal R^*\Psi}(\xi)|^2\,d\xi
    \le\frac{C_0\norm{\widehat\eta}_{L^1}}{(2\pi)^2}
    \int_{\Sph^{d-1}}\int_\R
    \langle\omega\rangle^{2r}|G(w,\omega)|^2\,d\omega\,d\sigma(w).
    \label{eq:localized-radon-final}
\end{align}
Since $2\bigl(r+\tfrac{d-1}2\bigr)=2r+d-1$, the left-hand side of
\eqref{eq:localized-radon-final} is a constant multiple of
$\norm{\eta\,\mathcal R^*\Psi}_{H^{r+(d-1)/2}(\R^d)}^2$, while by Plancherel in $b$
the right-hand side is a constant multiple of
$\norm{\Psi}_{L^2(\Sph^{d-1};H^r(\R))}^2$. Recalling $\Psi=\psi\Phi$ and the
multiplier bound $\norm{\psi\Phi}_{L^2(\Sph^{d-1};H^r(\R))}\le
C_\psi\norm{\Phi}_{L^2(\Sph^{d-1};H^r(\R))}$ proves
\eqref{eq:localized-dual-radon-estimate}.
\end{proof}

\end{document}
